\documentclass[11pt]{article}

\usepackage[margin=1in]{geometry}
\usepackage[T1]{fontenc}
\usepackage[utf8]{inputenc}
\usepackage{lmodern}
\usepackage{amsmath,amssymb}
\usepackage{booktabs}
\usepackage{array}
\usepackage{enumitem}
\usepackage{longtable}
\usepackage{xcolor}
\usepackage{hyperref}
\usepackage{listings}

\hypersetup{
    colorlinks=true,
    linkcolor=blue,
    urlcolor=blue
}

\lstset{
    basicstyle=\ttfamily\small,
    breaklines=true,
    columns=fullflexible,
    keepspaces=true,
    frame=single
}

\title{Mapping System: Classical vs.\ Cyclic}
\author{}
\date{}

\begin{document}
\maketitle

\section{Classical Mapping}

Classical mapping is intentionally simple. For every transition slot whose two endpoint vertices are both free, the slot is made bidirectional:
\begin{itemize}[leftmargin=2em]
    \item horizontal slots become left-and-right,
    \item vertical slots become up-and-down.
\end{itemize}

Conceptually, this is the ordinary undirected-grid interpretation expressed in explicit composite-grid form.

\section{Cyclic Mapping as a Pipeline}

Cyclic mapping is not just a raw directed overlay. The current implementation is a three-stage pipeline:
\begin{enumerate}[leftmargin=2em]
    \item overlay a deterministic cyclic direction pattern,
    \item clean up transitions that are invalid because of obstacle contact,
    \item postprocess the result to reduce unnecessary bidirectionals and restore connectivity where needed.
\end{enumerate}

\section{Raw Cyclic Overlay}

The raw cyclic pattern is encoded in \texttt{dev/maps/cycle\_pattern\_mapper.py}. It uses index parity and modulus classes such as \(i \bmod 4\) and \(j \bmod 4\) to orient transition slots. The effect is a repeating local circulation pattern over the composite grid.

The important point is not the exact modulus arithmetic but the design intent:
\begin{itemize}[leftmargin=2em]
    \item many transitions become one-way,
    \item local motion is given a preferred circulation direction,
    \item the resulting map is more structured than the classical fully bidirectional case.
\end{itemize}

\section{Obstacle-Aware Cyclic Cleanup}

The raw overlay alone is not safe enough. A directed transition can point into or out of an obstacle-adjacent area in a way that breaks the intended local cycle.

The cleanup stage therefore does two things:
\begin{enumerate}[leftmargin=2em]
    \item mark one-way cyclic transitions that touch obstacles,
    \item repair the surrounding local cycle by converting surviving related transitions into bidirectional transitions when needed.
\end{enumerate}

This repair step matters because the project does not want the raw cyclic stencil to create nonsense movement around blocked cells.

\section{Connectivity Postprocessing}

After obstacle cleanup, the cyclic map is passed through a second postprocessor:
\begin{itemize}[leftmargin=2em]
    \item first remove redundant bidirectional transitions if their endpoints remain mutually reachable through another route,
    \item then add bidirectional transitions back at empty slots if they are necessary to restore mutual reachability between adjacent free vertices.
\end{itemize}

The effect is a compromise:
\begin{itemize}[leftmargin=2em]
    \item keep the cyclic character where possible,
    \item but do not let the map fragment unnecessarily.
\end{itemize}

\section{Important Interpretation}

The current cyclic mapping is \textbf{not} a purely one-way map. After cleanup and connectivity restoration, some local edges may become bidirectional. This is deliberate. The project is comparing a strongly direction-biased map structure against the classical fully bidirectional structure, not forcing one-way motion at all costs.

\section{Why This Matters for Future Updates}

Any change to cyclic mapping should preserve the distinction between:
\begin{enumerate}[leftmargin=2em]
    \item the raw directed pattern,
    \item obstacle-aware repair,
    \item global connectivity repair.
\end{enumerate}

Collapsing these stages into one opaque step would make the project harder to reason about and would remove one of the clearest conceptual separations in the repository.

\end{document}
